\begingroup
\clearpage% Manually insert \clearpage
\let\clearpage\relax% Remove \clearpage functionality
\vspace*{-16pt}% Insert needed vertical retraction
\chapter[CONCLUSION]{CONCLUSION}
\endgroup

The LHC has entered the precision era. Discovering the Higgs boson was the beginning, not the end. The task now is to measure its couplings, probe the quantum structure of its interactions, and build the tools that will carry the program through the HL-LHC. This thesis contributed to each of these fronts, using modern graph and attention-based architectures to unlock physics that would otherwise remain buried in the complexity of collider events.

\begin{itemize}
    \item \textbf{Heavy-flavor jet tagging and calibration:} The transition from classical multivariate taggers to graph neural network and Transformer-based architectures (GN1, GN2, GN2X) delivered substantial improvements in b-jet and c-jet identification. The $t\bar{t}$ mistag calibration of GN2X in data validated the boosted tagger performance and provided the scale factors required to use it in ATLAS analyses. This work directly enables the Higgs measurements in Chapter~\ref{chap:higgs} and underpins the flavor tagging performance across the ATLAS Run~3 program.

    \item \textbf{Higgs boson Yukawa coupling measurements:} Two analyses targeting hadronic Higgs decays advanced the measurement of the Higgs Yukawa sector. In the $VH$ channel, the addition of a new low-$p_\mathrm{T}$ signal region achieved the first $5\sigma$ observation of $Z \to c\bar{c}$ and improved the upper limit on $H \to c\bar{c}$ by $\sim$4\%. In the VBF channel, a full analysis from sample production through statistical combination reported the first $3\sigma$ evidence for $VBF\,H \to b\bar{b}$ and the first upper limit on $VBF\,H \to c\bar{c}$. These results represent direct measurements of the Higgs couplings to second- and third-generation quarks.

    \item \textbf{Pileup mitigation for the HL-LHC:} At $\langle\mu\rangle = 200$, traditional pileup rejection algorithms degrade severely. PUMiNet and PhyGHT, an event-wide attention model and a hypergraph Transformer respectively, were shown to recover the Higgs boson, W boson, and top quark mass peaks under HL-LHC pileup conditions. An open-source simulated dataset was released alongside this work to enable reproducible benchmarking by the broader machine learning community. This study provides a foundation for the pileup mitigation strategies that will be needed when the HL-LHC begins operations.

    \item \textbf{Hadronic top quark polarimetry:} A proof-of-concept study demonstrated that top quark spin information can be extracted in the semi-leptonic $t\bar{t}$ channel without polarization-dependent training data. A Transformer regression model trained on unpolarized $t\bar{t}$ was evaluated on left- and right-handed polarized samples, recovering a measurable L/R asymmetry in $\cos\theta$ for both the bottom quark and the down quark. A cosine similarity-based purity cut was shown to recover the analyzing power $\beta$ toward truth-level expectations. This establishes a new avenue for hadronic top polarimetry that, with full ATLAS simulation and real LHC data, could contribute to measurements of $t\bar{t}$ spin correlations and quantum entanglement.
\end{itemize}

Taken together, these contributions demonstrate that modern attention-based architectures do more than improve benchmarks, they open experimental access to physics that was previously out of reach.
The LHC program is entering an era in which the bottleneck is no longer luminosity but the ability to fully exploit the information already present in the data.
The work presented here is a step toward that goal.
